\scanpage{46}\mspage{45}
de définir une figure $F'$ comme \unsure{un} \uline{\unsure{raffinant|raffinement}} de $F$, si chaque $X'\in F'$ \deleted{\unsure{déf}} est un raffinement d'un $X\in F$. Donc il vaut mieux prendre comme notion \unsure{primitive|première} celle de \uline{subdivision} d'une \deleted{\unsure{mul}} figure, (qui \unsure{va|en} se \unsure{ramener|ramène} à celle d'une multistrate), et définir un raffinement d'une figure comme sous-figure d'une subdivision.

Donc, nouveau départ !

Une \unsure{contrée} $C$ est définie par un triple
\[
 (\unsure{\mathcal M|M},\unsure{\mathcal F|\mathscr F},\Sigma)
\]
où $\mathcal M$ est un \unsure{ens} (ens. des \uline{multistrates}), $\mathcal F$ un ensemble de parties de $\mathcal M$ ($\mathcal F\unsure{\in|\subset}\unsure{\mathscr P|\mathfrak P}(\mathscr P(\mathcal M))$), \inserted{\unsure{appelées|appelés} \uline{figures},} et $\Sigma$ une relation sur $\mathcal F$ dite : $F'$ \uline{subdivise} $F$. \unsure{Axiomes|Axiome}
% slanted
\margin{Si $F\in\mathcal F$, les $X\in F$ sont les (multi)strates de $F$}

\uline{C 1} Pour tout $X\in\mathcal M$, $\exists$ \unsure{une|un} plus petit\unsure{e|$\subset$} \deleted{partie} figure \deleted{\unsure{qui}} $F_X\unsure{\subset|{,}}$ telle que $X\in F$.

On l'appelle \unsure{la} \uline{figure élémentaire} définie par $X$. On écrit
\[
 X\leq Y\ \overset{\mathrm{def}}{\Longleftrightarrow}\ F_X\subset F_Y\qquad\text{i.e. }X\in F_Y
\]
et on trouve une relation d'\uline{ordre}$^{*}$ $\leq$ sur $\mathcal M$.\textsuperscript{$*$}
% Two vertical strokes run beside C 1.
% slanted
\margin{$^{*}$ On veut que ce soit une relation d'ordre et pas seulement de préordre, i.e. \uline{C$'$1} si \deleted{\unsure{pr}} $X\ne Y$, $\exists$ $F\in\mathcal F$ qui \unsure{contient} un des deux \unsure{ms} $X,Y$, et pas l'autre, i.e. $F_X\ne F_Y$.}
% slanted
\margin{\unsure{--} i.e. $\mathcal F$ sépare les éléments de $\mathcal M$}
